\documentclass[french]{book}
\usepackage{teach}
\usepackage{amsmath,amsfonts,amssymb}
\usepackage{pgfplots}
%\usepackage{geometry}
%\geometry{top=1cm, bottom=1cm, left=2cm, right=2cm}
\usepackage[utf8]{inputenc}

%\usepackage[latin1]{inputenc}
%\usepackage[T1]{fontenc}

\usepackage[french]{babel}
\redefinecolor{thm}{title}{bgcolor}{alizarin}
\redefinecolor{prop}{title}{bgcolor}{meatbrown}
\redefinecolor{defn}{title}{bgcolor}{olivine}
\definecolor{alizarin}{rgb}{0.82, 0.1, 0.26}
\definecolor{meatbrown}{rgb}{0.9, 0.72, 0.23}
\definecolor{olivine}{rgb}{0.6, 0.73, 0.45}
\usepackage{pgf,tikz}

\usepackage{mathrsfs}
\usetikzlibrary{arrows}
\usetikzlibrary{patterns}

\begin{document}
\thispagestyle{empty}

\section{Devoir à la maison (niveau I)}
1) Rappeler l'aire d'un disque, notée $A_D$, en fonction de son rayon $R$.

2) Déterminer à quel intervalle appartient le rayon du cercle $R$ sachant que l'aire du disque soit plus grande ou égale à  $6\pi$ et strictement plus petite à $9\pi$.\\


\textit{On accepte et on utilisera pour la question 2) le résultat suivant: soit $0 \leq a \leq x \leq b$ alors $\sqrt{a} \leq \sqrt{x} \leq \sqrt{b}$ }

\section{Devoir à la maison (niveau II)}

\subsection*{La légende de Didon et le problème isopérimètrique}

\textit{La légende raconte qu'au IXème siècle avant Jésus-Christ, Elissa, princesse de Tyr (en Phénicie, actuels Liban, Israël et Syrie) devient reine de Tyr. Mais son frère Pygmalion assassine son époux afin de prendre le pouvoir. Horrifiée, Elissa s'enfuit. Elle atteint l'Afrique du Nord à Byrsa ("la peau de bœuf"), citadelle proche de l'actuel Tunis, et demande asile aux autochtones. On ne lui concède que ce que pourrait couvrir la peau d'un bœuf. Astucieuse, elle découpe la peau en si fines lanières qu'elle obtient, bout à bout, une longueur fantastique: près de 4 km. Avec la corde ainsi formée, elle aurait encerclé son territoire et fondé (vers 814 av. J.-C.) la très célèbre ville de Carthage actuellement nommée Istanbul.}\\



L'idée de former un cercle plutôt qu'un triangle, un rectangle, un carré ou tout autre forme géométrique fermée, place Didon au pinacle des mathématiques : elle avait donc admis sans hésiter le résultat isopérimétrique ci-après que Jacques Bernoulli prouva dans le cadre du calcul des variations (en 1698):



\begin{center}

\underline{ \textit{"De toutes les courbes fermées, de longueur donnée, celle qui entoure l'aire la plus grande est le cercle."} }
\end{center}

\section*{Questions}

1) Rappeler le périmètre $P$ d'un disque en fonction de son  rayon $R$, en déduire le rayon $R$ d'un disque en fonction du périmètre $P$ puis donner l'aire d'un disque en fonction de son périmètre. \\

2) Rappeler le périmètre $P$ d'un carré en fonction de son coté $c$, en déduire le coté $c$ d'un carré en fonction du périmètre $P$, puis donner l'aire d'un carré en fonction de son périmètre. \\

3) Montrer que, à périmètre fixé, l'aire d'un disque est toujours plus grande que l'aire d'un carré.

\section*{Une dimension plus loin... le problème isosurfacique}


\begin{center}
\textit{+2 points bonus au prochain DS en cas de réussite}
\end{center}

Pour les plus motivés et aguerris d'entre vous, voici un travail facultatif de recherche, notamment sur la racines cubiques. On appelle racine cubique d'un nombre réel $y$, l'unique solution $x$ de l'équation : $x^{3}=y$, on note cette solution $x={\sqrt[ 3]{y}}$. Cette notion sera nécéssaire pour les prochaines questions.\\

1f) Rappeler la surface $S$ d'une sphère en fonction de son rayon $R$, en déduire le rayon $R$ d'une sphère en fonction de la surface $S$ puis donner le volume d'une sphère en fonction de sa surface. \\

2f) Rappeler la surface $S$ d'un cube en fonction de son coté $c$, en déduire le coté $c$ d'un carré en fonction de la surface $S$, puis donner le volume d'un cube en fonction de sa surface. \\

3f) Montrer que, à surface fixé, le volume d'une sphère est toujours plus grande que le volume d'un cube.\\




\end{document}